\section{Problema y alcance del avance}
Las fuentes urbanas publican datos con diferentes estructuras, frecuencias,
reglas de calidad e identificadores. SmartCity Analytics estudia cómo almacenar
estos datos y consultarlos sin perder su procedencia ni atribuirles información
que no existe en los originales.

El primer avance se concentra en PostgreSQL y MongoDB: modelos, restricciones,
índices, población de datos, diez consultas por motor y evidencia de ejecución.
El plan del equipo sitúa este avance el 28 de octubre y la entrega final el
25 de noviembre de 2026. Hadoop, Spark y la interfaz corresponden al hito final.
La división procede del plan de trabajo; el enunciado general no incluye una
rúbrica independiente del primer avance \cite{plan,enunciado}.

\subsection{Objetivos del avance}
\begin{itemize}
\item Representar catálogos y observaciones con integridad relacional.
\item Representar lecturas y sus metadatos como documentos BSON validados.
\item Demostrar consultas, agregaciones, índices y operaciones CRUD solicitadas.
\item Comprobar replay, conflictos y respaldo/restauración en entornos locales.
\item Documentar qué resultados están acreditados y qué requisitos faltan.
\end{itemize}

\section{Fuentes, preparación y procedencia}
Las fuentes consideradas son calidad del aire, meteorología y tráfico del
Ayuntamiento de Madrid. Sus catálogos oficiales y la investigación del repositorio
documentan licencia CC BY 4.0; deben conservarse atribución, snapshot y cambios
realizados \cite{aire,meteo,trafico,contrato}.

\begin{longtable}{p{0.18\textwidth}p{0.34\textwidth}p{0.38\textwidth}}
\toprule
Fuente & Información & Estado en la generación anual \\
\midrule\endhead
Aire & Observaciones horarias y flags de calidad. &
1.076.813 observaciones aceptadas de 2025. \\
Meteorología & Variables y estaciones meteorológicas. &
Los 69.408 slots auditados de enero permanecen en cuarentena;
el reloj no está acreditado para habilitar su aceptación. \\
Tráfico & Intervalos y variables publicadas por punto de medida. &
677.724 filas seleccionadas de enero permanecen en cuarentena;
la revisión incompatible del contrato sigue pendiente. \\
\bottomrule
\end{longtable}

Solo aire constituye una fuente de mediciones aceptada en la carga anual.
Investigar tres fuentes o disponer de inventarios auxiliares no demuestra por sí
solo el requisito de dos fuentes integradas.

\subsection{Generación local v2: dos fuentes, sin modificar v1}
La implementación local autorizada el 27/10 prepara un paquete nuevo de
\textbf{748.889 lecturas reales de enero de 2025}: 91.959 de aire y 656.930 de
tráfico. Se excluyen 17.774 registros C30 y 3.020 flags vacíos de tráfico, además
de los rechazos de aire. No se suman estos datos a la carga anual v1 ni se
presentan como una ampliación de 1.733.743 registros.

\archivo{sprint1/basico/preparar_v2.py} crea la generación
\archivo{expanded-v2-2025-01} con cuarentena y manifiesto de hashes.
La preparación está terminada; las cargas nativas y las consultas deben
acreditarse con el estado y las salidas de los nuevos informes de ejecución.
Disponer del archivo preparado no demuestra por sí solo una carga en cada motor.

Como complemento separado se preparan 14.388 accidentes de 2026,
deduplicados desde 33.322 filas de personas \cite{accidentes}.
Se eliminan sexo, edad, lesiones, dirección y expediente en claro. Las fechas
se conservan como fechas publicadas sin inventar un reloj UTC; no se unen
estos accidentes de 2026 con las observaciones de enero de 2025.

\subsection{Calidad y conteos}
El snapshot de aire produjo 1.103.232 slots de 2025: 1.076.813 aceptados y
26.419 rechazados. Los rechazos corresponden a 26.167 flags del publicador,
126 horas locales inexistentes y 126 ambiguas sin evidencia para resolverlas.
La ampliación anual contiene enero; no se le suman de nuevo sus 91.959 aceptadas
ni las muestras de prueba \cite{cierre}.

El campo de tráfico \archivo{periodo_integracion} cuenta muestras, no minutos.
La duración documentada es de 15 minutos. Esta corrección está en la propuesta
v2 y no habilita automáticamente la carga de tráfico bajo v1 \cite{v2}.

La versión local v2 aplica esa corrección en una base nueva: duración 15 y
conteo de muestras positivo, incluso cuando el conteo es 30. Solo admite URB/M30
con flag N y tiempo local inequívoco. Los negativos opcionales se marcan como
faltantes; URB no recibe una velocidad inventada y carga es un índice.
Los distritos provienen del inventario disponible, sin vigencia histórica
confirmada para enero de 2025.

\subsection{Identidad y trazabilidad}
Se distinguen estación, canal lógico y sensor físico. Un canal no acredita la
identidad de un dispositivo individual. El contrato anual deriva identificadores
mediante SHA-256 de contenido canónico y conserva fuente, versión del archivo,
referencia de fila/columna, calidad, tiempo original y transformación.

\archivo{source_version} identifica el snapshot y debe coincidir con
\archivo{raw_reference.sha256}. Las horas de aire se interpretan en
\archivo{Europe/Madrid}; las ambigüedades no se resuelven inventando un fold.
Los conflictos de contenido no deben sobrescribir observaciones aceptadas.

\subsection{Datos administrativos sintéticos}
La evidencia documental de la respuesta docente del 10/10 permite atributos
administrativos aleatorios. El paquete anual registra una organización,
5.133 usuarios y 5.133 sensores sintéticos deterministas, con semilla
\archivo{20261010}, versión \archivo{admin-demo-v2} y procedencia explícita.
Los correos usan \archivo{example.invalid}. Esto no acredita inventario físico
ni autoriza mediciones inventadas \cite{permiso,contrato}.

Los sensores demo tienen identificadores propios, pero no se enlazan a canales
o lecturas observadas: la referencia \archivo{sensor_id} en los canales y en
la proyección documental de las lecturas reales permanece nula.
La autorización se conoce por el registro
aportado al repositorio, no por una captura original adjunta a este informe.
